% ch3 B.2 后半 (原书页 119–125 / PDF p134–p140)

\begin{equation*}
E_{\Delta}(k_{0})\;\, \Delta(r)\;\, q_{k_{0}}^{\dagger}(r)\;\, \Psi_{g}
\end{equation*}

与式 (8) 相符。

于是，在式 (9) 中，$E(k)$ 扮演的正是式 (7) 中散射矩阵元 $M_{k,\mathbf{K}}$ 在波长极长的
声子极限（$\mathbf{K}\to 0$）下所扮演的角色，只不过 $M_{k,\mathbf{K}}$ 乘的是位移 $Q$
而不是形变 $\Delta$。虽然用 Hartree-Fock 理论作直接计算来证明 $M_{k,\mathbf{K}}$ 正比于
$\mathbf{K}$ 并不困难，但我们给出的这一论证在存在大的多体效应时依然有效，并且把散射矩阵元
直接同实验上可观察的现象联系起来。

鲜为人知的是，上述论证在那些对称性允许压电效应（piezoelectric effect）的晶体中会完全
失效。在这类晶体中，形变——特别是切变形变 $Y$——不仅能改变电子的有效势，还能建立起一个
正比于应变的电场 $E$。这就给我们一个势项

\begin{equation*}
E'(r) \;-\; E(r) \;\propto\; V(r) \;=\; e\mathbf{E}\cdot\mathbf{r} \;\propto\; rS(r)
\end{equation*}

它正比于总位移而不是应变。因此在这种情形下，所谓"形变势定理"并不成立。Hopfield 曾分析过
由此可能引出的若干相当奇怪的后果；它会导致慢电子极为反常的散射与极化效应。

然而，在较常见的金属和半导体中，这一定理确实成立。这一点在两条或更多条能带发生简并或
近简并的情形下尤为重要：定理此时仍然有效，但最方便的做法是让形变势 $E$ 推广为能带指标下的
一个一般矩阵 $E^{nn'}$。在这种情形下，要从刚性离子模型直接算出正确答案就不那么容易了。

在这种电子-声子相互作用存在时，准粒子 $q_{k}^{\dagger}$ 可以在布里渊区（Brillouin zone）
内从一点被散射到另一点。在绝对零度、不存在声子时，能带中绝对最低的准粒子态
$q_{k}^{\dagger}\Psi_{g}$ 仍然不会被散射，因为它无法发射声子——它是具有给定动量、能量和
粒子数的唯一状态；但动量较高的准粒子态一旦满足
$\mathrm{d}E/\mathrm{d}k > \hbar(\mathrm{d}\omega/\mathrm{d}k)_{\text{phonon}} = \hbar c$，
就能够伴随一个声子的发射而被散射，其中 $c$ 是材料中最低的声速（图 35）。若准粒子的有效质量
为 1，而 $c$ 取典型值 $10^{5}$ cm/sec，这就是

\begin{align*}
\frac{\hbar^{2}}{m}\,(k - k_{0}) &> \hbar c \\
k - k_{0} &> \frac{mc}{\hbar} \;\sim\; 10^{5}\ \text{cm}^{-1} \\
\frac{\hbar^{2}}{2m}\,(k - k_{0})^{2} &\geq \frac{mc^{2}}{2} \;=\; \frac{1}{2}\,10^{-17}\ \text{erg} \;\approx\; 0.04\,^{\circ}\mathrm{K}
\end{align*}

这是一个低得惊人的能量，因此从实用角度看，我们从来不会遇到不能通过发射声子而被散射的
准粒子。

%FIG{35}: 图 35

我们立刻意识到，这迫使我们修改准粒子元激发的原始定义，因为在这种情形下，准粒子态不再是
整个系统的严格本征态。这是一种普遍现象：在几乎所有的情形中，由于总存在这种或那种散射的
可能性，元激发态都不是完全严格的本征态。

就我所知，还从未有人在有散射存在时给出过准粒子的绝对严格的数学定义；但定义准粒子的方式
有若干种，在任何一个给定的情形中它们都可以是令人满意的。这些方式可以归入两条思想路线，
我们分别称之为"赝哈密顿量"（pseudo-hamiltonian）方法和"完全重整化"（full
renormalization）方法。

这两条处理问题的一般思想路线中的第一条，本质上就是我们在讨论形变势时已经实行过的那套
做法。也就是说，我们尽最大努力把准粒子与晶格运动彼此完全独立地定义出来，并把它们的相互
作用显式地留在问题之中。在形变势的讨论中，我们实际上是走进固体内部，给所有离子实装上刚性
夹钳，把它们绝对固定在与某种长波长晶格形变相应的位置上。由此我们得到准粒子能量对形变的
某种依赖关系，而当我们松开夹钳时，它就充当了电子-声子相互作用能。用形式化的语言说，我们
为准粒子建立了一个有效哈密顿量

\begin{equation*}
\mathcal{H}_{\text{eff}} \;=\; \int \mathrm{d}r \int \mathrm{d}r'\;\, q^{\dagger}(r)\, H_{\text{eff}}(r,r')\, q(r')
\end{equation*}

其中 $H_{\text{eff}}(r,r')$ 是局域形变 $\Delta(r)$ 和 $S(r)$ 的线性函数。于是我们可以把
$\Delta(r)$ 和 $S(r)$ 本身作为多体问题的量子化坐标引入，它们将带来一个必须添入的、二次型
的自能 $\sum_{k}\hbar\omega_{\mathbf{K}}^{0}\,(P_{\mathbf{K}}^{2} + Q_{\mathbf{K}}^{2})$。
我们希望，相互作用项随后可以由微扰论计入，并导致准粒子的散射，如下：散射哈密顿量就是
形变势，即式 (9)：

\begin{equation*}
\mathcal{H}' \;=\; E_{\Delta} \sum_{k\mathbf{K}} \Delta(\mathbf{K})\;\, q_{k+\mathbf{K}}^{\dagger} q_{k}
\end{equation*}

$\Delta(\mathbf{K})$ 与声子产生算符 $b_{\mathbf{K}}^{*}$ 之间关系的归一化，可以由能量
表达式得到：

\begin{equation*}
\frac{S}{2}\,\bigl|\Delta\bigr|^{2} \;=\; \text{势能} \;=\; \frac{1}{2}\,(\text{总能量}) \;=\; \frac{1}{2}\,\hbar\omega_{\mathbf{K}}\Bigl(n_{\mathbf{K}} + \frac{1}{2}\Bigr)
\end{equation*}

其中 $S$ 是压缩刚度 $\rho c^{2}$，而 $\omega_{\mathbf{K}} = c\mathbf{K}$ 是声子的频率。
于是我们可以取

\begin{equation*}
\Delta(\mathbf{K}) \;=\; \frac{1}{2}\sqrt{\frac{\hbar\mathbf{K}}{\rho c}}\;\bigl(b_{\mathbf{K}}^{*} + b_{-\mathbf{K}}\bigr)
\end{equation*}

由含时微扰论给出的、从单一态进入一个态连续统的逆平均自由时间（即跃迁速率）的表达式——
也就是从 $k$ 进入已自发发射了一个声子 $\mathbf{K}$ 的那些态的速率——是

\begin{equation*}
\frac{\hbar}{\tau} \;=\; 2\pi\,\bigl|\text{M.E.}\bigr|^{2}\,\frac{\mathrm{d}n}{\mathrm{d}E}
\end{equation*}

当然，M.E. 指的是散射矩阵元。物理上唯一有趣的情形，是电子动量相当大的情形，此时

\begin{equation*}
\frac{\hbar^{2}k^{2}}{2m} \;\gg\; \hbar k c\,,\qquad k \;\gg\; K_{0} = \frac{mc}{\hbar}
\end{equation*}

在这种情形下，声子造成的动量变化远大于能量变化，于是 $|k+\mathbf{K}| \cong |k|$
（图 36）。这时 $\mathrm{d}n/\mathrm{d}E$ 就恰好是电子谱中的态密度（density of states）。

%FIG{36}: 图 36

散射前的寿命如何计算，我留作练习：结果是

\begin{equation*}
\frac{\hbar}{\tau} \;=\; \frac{4\pi^{2}}{3}\, E^{2} \left(\frac{m^{2}}{\hbar^{3}\rho c}\right) \left(\frac{\hbar^{2}k^{2}}{2m}\right)
\tag{10}
\end{equation*}

若晶格常数为 $a$，则 $\rho \cong M/a^{3}$，而 $\hbar^{2}/ma^{2}$ 与电子能量 $\sim 1$ Ry
同量级；$\hbar c/a$ 与 $\sqrt{m/M} \times 1$ Ry 同量级。形变势常数 $E$ 同样可以预期为
1 Ry 的量级，因为百分之百的形变必然会使电子的能量完全改变；于是我们可以说

\begin{equation*}
\left.\frac{\hbar}{\tau}\;\right|_{\;\left(E_{k} \,=\, \frac{\hbar^{2}k^{2}}{2m}\right)} \;\sim\; \sqrt{\frac{m}{M}}
\end{equation*}

我们看到 $\hbar/\tau$ 确实非常小，而且对能量最低的准粒子令人满意地趋于零。

与散射相伴随的，是在等能面附近有一些态被混杂进来：对给定的能量差 $\Delta E$，微扰论给出
这部分混杂为

\begin{equation*}
\frac{\bigl|\text{M.E.}\bigr|^{2}}{(\Delta E)^{2}}\;\,\mathrm{d}n
\end{equation*}

于是，在能量上离开准粒子态超过 $\Delta E$ 的那些态的总混杂为

\begin{equation*}
\bigl|\text{M.E.}\bigr|^{2} \int_{\Delta E}^{\infty} \frac{\mathrm{d}n}{\mathrm{d}E}\,\frac{\mathrm{d}E}{E^{2}} \;\cong\; \frac{\hbar}{\tau}\,\frac{1}{\Delta E}
\end{equation*}

显然，如果允许任意低能量的态混入，此式就会发散。这意味着，任何想在有散射存在时把准粒子
态定义为严格本征态的尝试都注定失败。正如我所强调过的，我们可以沿两条途径处理这个问题：
第一条就是我们刚才描述的有效哈密顿量理论，其中准粒子是一个"被夹住的"（clamped）系统的
严格本征态，而电子-声子相互作用保留在总哈密顿量之中，我们试图用这个总哈密顿量去求解任何
输运问题等等。

与之相对的第二种方法，是试图把系统当作一个整体来处理，并定义由于散射而随时间衰变的准粒子
态。这是现代多体理论家更青睐的方法；见 Nozières (68)。正如我们稍后将要讨论的，这样的
准粒子态不仅带有重整化常数 $Z$，还带有复能量 $E_{k} + i(\hbar/\tau)$；我们说它们是
"格林函数（Green's function）的复极点"。它们与完全重整化理论的传播子（propagator）
相联系；有效的电子-声子相互作用可以作为理论的重整化顶角（vertex）部分引入，等等。要在
这种完全重整化的形式下开展整个输运理论等等，目前还差不多只是一个纲领，但毫无疑问它完全是
可行的。在半导体和绝缘体中，第一种方法迄今为止肯定是足够的，而且它免去了随身携带诸如 $Z$
这类各种非物理重整化参数的必要。

尽管如此，也有人可以争辩说第一种方法有点含糊而且不合乎物理，因为真实的准粒子随身携带着
一大片声子极化云，尤其是短波长声子和光学声子的极化云。但我们可以设想，只对那些真正能使
准粒子受到散射的长波长声子加上夹钳；这几乎不会影响准粒子的极化云，却使我们能够相当严格地
定义它们。特别是，准粒子既然与 $N+1$ 体系统的严格本征态相关联，就具有固定的动量并携带
电荷 $e$，只是处理 $Q\to 0$ 极限时须加小心——这一点我们已经讨论过了。在存在场时，它们将
服从通常类型的输运方程，带有通常的 Lorentz 力。

这种处理绝缘体中准粒子问题的方法看起来是完全严格而准确的，至少对非压电晶体和足够低的
温度是如此。然而，在电子-声子耦合极强的某些情形中，伴随电子行进的极化云的效应可能相当大。
这在某些 d 带氧化物晶体中是典型的：给定离子上 d 壳层能级中的一个电子可以使周围的离子发生
相当严重的极化，其幅度相对于这些离子的零点涨落而言相当可观。电子跳跃到另一个 d 壳层离子
上去的矩阵元的幅度可以写成

\begin{align*}
b_{d_{1}d_{2}} &= \bigl(\Psi_{d_{1}}\,|\,\mathcal{H}\,|\,\Psi_{d_{2}}\bigr) \\
&\cong \bigl(a_{d_{1}}\,|\,\mathcal{H}_{\text{el}}\,|\,a_{d_{2}}\bigr) \bigl(\Psi_{\text{lattice}}(\text{电子在 } d_{1}\text{ 上})\,\big|\,\Psi_{\text{lattice}}(\text{电子在 } d_{2}\text{ 上})\bigr)
\end{align*}

$a$ 是 d 带的普通 Wannier 函数，而 $\Psi_{d}$ 是完整的多体波函数。第二个因子——晶格
波函数的交叠——将会很小，这就大大增加了准粒子的有效质量。这就是与"自陷极化子"相对应的
物理图像——它能够运动，但只能带着大为减小的矩阵元和大为增大的 $m^{*}$ 运动。甚至在相当
低的温度下，电子通过热激活而不是通过这种隧穿过程来运动的概率也可能变得很大，这意味着这类
材料中的输运根本不像普通半导体中的输运（文献 (71)）。

在更高温度下，即使在正常情形中，热激发声子造成的额外散射也会变得远大于发射所致的自发
散射。声子寿命中的（矩阵元）$^{2}$ 因子必须再乘上一个占据因子：发射时为 $n+1$，吸收时为
$n$，其中 $n = 1/[\exp(\hbar\omega/k_{\mathrm{B}}T) - 1]$。对半导体中那些有关的声子——
即满足

\begin{equation*}
\mathbf{K}_{\text{phonon}} \;\sim\; k_{\text{electron}}
\end{equation*}

的声子——这个因子可以变得相当大，其中

\begin{equation*}
k_{\text{electron}} \;=\; \frac{\sqrt{2m^{*}k_{\mathrm{B}}T}}{\hbar}
\end{equation*}

\begin{equation*}
n(\mathbf{K}) \;\sim\; \frac{k_{\mathrm{B}}T}{\hbar c\mathbf{K}} \;=\; \frac{V_{\text{el}}}{c} \;\sim\; 10 \qquad \text{（在 } 50{-}100\,^{\circ}\mathrm{K}\text{ 时）}
\end{equation*}

最慢的那些电子所受的热散射可以变得极端严重，即使在低温下也是如此，尽管这在物理上似乎并不
十分重要。在更高温度下，我们的散射率 $\hbar/\tau$ 公式 (10) 要乘以 $k_{\mathrm{B}}T/\hbar c\mathbf{K}$
（译注：原书此处印作 $k_{0}T$，按上下文应为 $k_{\mathrm{B}}T$），并且由于 $\mathbf{K}$
与 $k$ 同量级、$k_{\mathrm{B}}T$ 按 $k^{2}$ 变化，我们得到

\begin{equation*}
\frac{\hbar}{\tau} \;\sim\; \frac{E^{2}}{\rho c^{2}}\, k^{3} \;\sim\; T^{3/2}
\end{equation*}

迁移率 $\mu = e\tau/m$ 的 $T^{-3/2}$ 依赖关系是半导体中众所周知的结果。注意，随着 $T$
增大，$\hbar/\tau$ 将相当快地变得比 $k_{\mathrm{B}}T$ 还大。

我不知道有任何关于这种情形的讨论；在其中，若没有进一步的论证，准粒子概念是不应使用的。
它出现在室温下迁移率为 10 左右的区域——在劣质半导体中这是并不罕见的一类数值。

金属中电子-声子散射变强时的效应，比起绝缘体和半导体中的情形来，被人们理解得少得多。
金属中在低温下，费米面（Fermi surface）附近的电子不能发生自发发射，因为末态电子态已被
占满。导致电子散射寿命的热吸收与热发射过程有两种不同的温度依赖关系。低温下只有长波长声子
被激发。各位当会记得，散射矩阵元正比于 $\mathbf{K}^{1/2}$，而可能的末态近似地位于
$\mathbf{K}$ 空间中的一个球面上，其面元 $\propto \mathbf{K}\,\mathrm{d}\mathbf{K}$。于是

\begin{equation*}
\frac{\hbar}{\tau} \;\propto\; \int_{0}^{\mathbf{K}_{F}} \frac{\mathbf{K}^{2}\,\mathrm{d}\mathbf{K}}{\exp(\hbar c\mathbf{K}/k_{\mathrm{B}}T) - 1} \;\propto\; \left(\frac{k_{\mathrm{B}}T}{\hbar c}\right)^{3}
\end{equation*}

在高于 Debye 温度的温度下，所有声子态的占据数都正比于 $T$，而
$\hbar/\tau \propto T/\theta_{D}$（图 37）。

我们可以在高温极限下区分两种情形：$\hbar/\tau$ 大于或小于 $kT$。当 $\hbar/\tau < kT$
时，无论哪一类微扰论大概都会相当令人满意，

%FIG{37}: 图 37
